\documentclass[10pt]{article}
\usepackage[T1]{fontenc}
\usepackage[utf8]{inputenc}
\usepackage[french]{babel}
\usepackage{amsmath,amssymb,mathrsfs,mathpazo}
\usepackage{geometry}
\geometry{a4paper,margin=2cm}
\usepackage{graphicx,array,tabularx,tabularray,colortbl,multirow,multicol,booktabs,enumitem,pifont,ulem,fancybox,diagbox,cancel}
\usepackage{pstricks,pstricks-add,pst-plot,pst-tree,pst-node,pst-text,pst-func,pst-math,pst-bspline,pst-3dplot,pst-eucl,pst-solides3d}
\let\pscalculate\undefined
\usepackage{pst-bezier}
\usepackage{tikz,pgfplots,tkz-tab}
\usetikzlibrary{arrows,arrows.meta,patterns,calc,positioning,trees,decorations.markings,shapes.geometric,shapes.misc,shapes.arrows,intersections,decorations.pathmorphing,decorations.pathreplacing,backgrounds,shadows,babel}
\pgfplotsset{compat=1.18}
\usepgfplotslibrary{fillbetween}
\usepackage[np]{numprint}
\usepackage{listings,makecell,pas-tableur,twoopt}
\let\barre\undefined
\usepackage[pstricks]{tablvar}
\definecolor{darkgreen}{RGB}{0,100,0}
\providecommand{\red}{\color{red}}
\providecommand{\blue}{\color{blue}}
\providecommand{\green}{\color{green}}
\providecommand{\black}{\color{black}}
\providecommand{\white}{\color{white}}
\providecommand{\gray}{\color{gray}}
\providecommand{\cyan}{\color{cyan}}
\providecommand{\magenta}{\color{magenta}}
\providecommand{\yellow}{\color{yellow}}
\providecommand{\violet}{\color{violet}}
\providecommand{\orange}{\color{orange}}
\providecommand{\euro}{\texteuro}
\providecommand{\EUR}{\texteuro}
\providecommand{\ieme}{\textsuperscript{e}}
\providecommand{\degres}{\ensuremath{^\circ}}
\providecommand{\degree}{\ensuremath{^\circ}}
\providecommand{\vv}[1]{\overrightarrow{#1}}
\providecommand{\Coord}[1]{#1}
\providecommand{\decosix}{}
\providecommand{\starredbullet}{\ensuremath{\bullet}}
\pagestyle{empty}
\setlength{\parindent}{0pt}
\usepackage{adjustbox}\DeclareUnicodeCharacter{2236}{:}\DeclareUnicodeCharacter{2190}{\ensuremath{\leftarrow}}
\providecommand{\R}{\mathbb{R}}
\providecommand{\N}{\mathbb{N}}
\providecommand{\Z}{\mathbb{Z}}
\providecommand{\Q}{\mathbb{Q}}
\providecommand{\D}{\mathbb{D}}
\providecommand{\C}{\mathbb{C}}
\providecommand{\e}{{\mathrm{e}}}
\providecommand{\vect}[1]{\overrightarrow{#1}}
\providecommand{\vectt}[1]{\overrightarrow{\mathrm{#1}}}
\providecommand{\vvt}[1]{\overrightarrow{\mathrm{#1}}}
\providecommand{\np}[2][]{{#2~#1}}
\providecommand{\vv}[1]{{\overrightarrow{#1}}}
\renewcommand{\d}{{\mathrm{d}}}
\providecommand{\degres}{{^\circ}}
\providecommand{\degre}{{^\circ}}
\providecommand{\blue}{\color{blue}}
\providecommand{\red}{\color{red}}
\providecommand{\enskip}{\quad}
\providecommand{\llbracket}{\lbrack\!\lbrack}
\providecommand{\rrbracket}{\rbrack\!\rbrack}
\providecommand{\euro}{{\text{€}}}
\providecommand{\no}{{\text{n°}}}
\providecommand{\interval}[2]{{[#1;#2]}}
\providecommand{\intervFO}[2]{{[#1;#2[}}
\providecommand{\dotfill}{\ldots}

\providecommand{\ds}{\displaystyle}
\providecommand{\textding}[1]{#1}
\providecommand{\Oij}{\(\mathrm{O};\vec{\imath},\vec{\jmath}\)}
\providecommand{\Oijk}{\(\mathrm{O};\vec{\imath},\vec{\jmath},\vec{k}\)}
\providecommand{\Ouv}{\(\mathrm{O};\vec{u},\vec{v}\)}
\providecommand{\intervFF}[2]{\ensuremath{[#1;#2]}}
\renewcommand{\euro}{\ifmmode\text{\texteuro}\else\texteuro\fi}
\DeclareUnicodeCharacter{00B0}{\ifmmode^{\circ}\else\textdegree\fi}
\DeclareUnicodeCharacter{20AC}{\ifmmode\text{\texteuro}\else\texteuro\fi}
% Énoncés : Éducation nationale. Transcription et corrigés : APMEP, Denis Vergès et contributeurs. Découpage et classement : CoopMaths, Sébastien Lozano, Éric Elter et contributeurs. Adaptation de la présentation : Maths Explorer.
% CC BY-SA : https://coopmaths.fr/www/about/
% https://coopmaths.fr/alea/static/bac/2023/tex/bac_2023_03_sujet2_etrangers_groupe1_3.tex
\begin{document}
\noindent\textit{APMEP / CoopMaths. CC BY-SA. Présentation adaptée par Maths Explorer.}\par\medskip
\medskip

La figure ci-dessous correspond à la maquette d'un projet architectural.

Il s'agit d'une maison de forme cubique (ABCDEFGH) accolée à un garage de forme cubique (BIJKLMNO) où L est le milieu du segment [BF] et K est le milieu du segment [BC].

Le garage est surmonté d'un toit de forme pyramidale (LMNOP) de base carrée LMNO et de sommet P positionné sur la façade de la maison.

\begin{center}
\psset{unit=1cm,arrowsize=2pt 4}
\begin{pspicture}(0,-4.5)(6.8,4.7)

\pspolygon[fillstyle=solid,fillcolor=lightgray](3.1,0.3)(2.2,-0.6)(3.4,-1.8)(4.85,-0.9)(3.65,0.3)
\psline[linewidth=1.25pt]{->}(0,0)(1,-1)\uput[d](0.5,-0.5){$\vect{\imath}$}
\psline[linewidth=1.25pt]{->}(0,0)(0,1.3)\uput[l](0,0.65){$\vect{k}$}
\psline[linewidth=1.25pt,linestyle=dashed]{->}(0,0)(1.5,1.1)\uput[u](0.75,0.55){$\vect{\jmath}$}
\psline(0,0)(4.2,-4.2)
\psline(0,1.1)(0,3.6)
\psline(0,2.6)(2.2,0.4)(2.2,-2.2)
\psline(0,2.6)(2.9,4.4)(5.1,2.2)(2.2,0.4)
\psline(5.1,2.2)(5.1,-0.4)(4.5,-0.8)
\psline[linestyle=dashed](2.2,-2.2)(5.1,-0.4)
\psline[linestyle=dashed](4.2,3)(1.5,1.1)
\psline(2.2,-0.6)(3.4,-1.8)(3.4,-3.4)
\psline(3.4,-1.8)(4.85,-0.9)(4.85,-2.5)(3.4,-3.4)
\psline(5.8,4.1)(4.22,3.)
\psline[linestyle=dashed](4.85,-2.5)(3.65,-1.3)(3.65,0.3)
\psline[linestyle=dashed](2.9,4.4)(2.9,2.1)(5.1,-0.4)
\psline[linestyle=dashed](3.65,0.3)(2.2,-0.6)
\pspolygon(3.1,0.3)(2.2,-0.6)(3.4,-1.8)
\pspolygon(3.1,0.3)(3.65,0.3)(4.85,-0.9)
\uput[l](0,0){A} \uput[dl](2.2,-2.2){B} \uput[r](5.1,-0.4){C} \uput[l](2.9,2.1){D}
\uput[l](0,2.6){E} \uput[u](2.2,0.4){F} \uput[ur](5.1,2.2){G} \uput[u](2.9,4.4){H}
\uput[dl](3.4,-3.4){I} \uput[dr](4.85,-2.5){J} \uput[r](3.65,-1.3){K} \uput[l](2.2,-0.6){L}
\uput[dl](3.4,-1.8){M} \uput[r](4.85,-0.9){N} \uput[ur](3.65,0.3){O} \uput[u](3.1,0.3){P}
\end{pspicture}
\end{center}

On munit l'espace du repère orthonormé $\left(\text{A}~;~\vect{\imath},~\vect{\jmath},~\vect{k}\right)$, avec $\vect{\imath} = \dfrac12\vect{\text{AB}},~\vect{\jmath} = \dfrac12\vect{\text{AD}}$ et $\vect{k} = \dfrac12\vect{\text{AE}}$.

\medskip

\begin{enumerate}
\item 
	\begin{enumerate}
		\item Par lecture graphique, donner les coordonnées des points H, M et N.
		\item Déterminer une représentation paramétrique de la droite (HM).
	\end{enumerate}
\item L'architecte place le point P à l'intersection de la droite (HM) et du plan (BCF).

Montrer que les coordonnées de P sont $\left(2~;~\dfrac23~;~\dfrac43\right)$.
\item
	\begin{enumerate}
		\item Calculer le produit scalaire $\vect{\text{PM}} \cdot \vect{\text{PN}}$.
		\item Calculer la distance PM.
		
On admet que la distance PN est égale à $\dfrac{\sqrt{11}}{3}$.
		\item Pour satisfaire à des contraintes techniques, le toit ne peut être construit que si l'angle $\widehat{\text{MPN}}$ ne dépasse pas $55\degres$.
		
Le toit pourra-t-il être construit ?
	\end{enumerate}
\item Justifier que les droites (HM) et (EN) sont sécantes.

Quel est leur point d'intersection ?
\end{enumerate}

\bigskip

\end{document}
